\scanpage{55}\mspage{6}
donner une section de $(X\smallsetminus\Delta)/\underline V$ \deleted{\unsure{de}}, en prenant \inserted{pour $I$} l'image inverse de cette section dans $X\smallsetminus\Delta$.

Mais la donnée d'une structure octaédrale équivaut à celle d'un \deleted{$\unsure{\underline I|\underline J}$} rev. étale de degré quatre de $S$ (à ne pas confondre avec le $\underline I$ de tantôt, appelons le $\underline K$), ou encore, à celle de

a) un schéma en groupes $V$ sur $S$, loc. isom. à $(\mathbb D_2)_S$ -- ou encore un revêtement fini étale $\underline J$ ($=\underline V^*$) de degré 3 de $S$, ou \deleted{de} enfin, à celui d'un \deleted{\unsure{gr}} fibré en droites projectives $\Sigma=\Sigma_{\underline J}$ sur $S$, muni d'un diviseur \inserted{relatif} étale $D_J$ de degré 3 (\deleted{\unsure{si}} isom. \deleted{\unsure{lui}} \inserted{à} $\underline J$)

b) d'un torseur $\underline K$ sous $\underline V$, ou ce qui revient au $=$, d'un revêtement \deleted{\unsure{galoisien}} \inserted{principal} \inserted{$X^*$} de $\Sigma_{\underline J}^*=\Sigma_{\underline J}\smallsetminus\deleted{\unsure{J}}D_{\underline J}$ de groupe $\underline V$, \deleted{\unsure{géométriquement connexe sur chaque fibre}} correspondant à \deleted{$\unsure{\xi}$} un élément de \deleted{$\unsure{H^0(S,R^1 f_* V)|H^0(S,R^1 f_*(V)}$} $H^0(S,\mathcal H^1(\Sigma_{\underline J}^*/S,\underline V))\simeq\mathrm{Hom}(\underline V_{\underline J},\underline V)$ \deleted{$\unsure{=\underline V_{\underline J}|/\mathrm{iso}}$} qui est l'identité $\underline V_{\underline J}\longrightarrow \underline V=\underline V_{\underline J}$.
% slanted
\margin{\uline{NB} Le syst. local des $(\pi_1)_{\mathrm{ab}/2}$ de $\Sigma_{\underline J}^*$ sur $S$ est \deleted{$\unsure{V}$} can. isom. à $V_{\underline J}$ \ldots}
-- La correspondance entre $\underline K$ et $X$ étant obtenue \deleted{par} la formule
\[
 X^*=\deleted{\unsure{X_{J,0}\wedge_{\Gamma}|\Sigma_{\underline J,0}\wedge_{\Gamma}}}\ K\mathbin{\wedge}_{V}\unsure{\overline\Sigma|\Sigma}_{\underline J,2}^*, \tag*{(11)}
\]
où $\overline\Sigma_{\underline J,2}^*$ est le revêtement principal \unsure{dié|dû|dé|dit}
% archive sheet number: 54
